\documentclass[11pt,a4paper]{article}
\usepackage[utf8]{inputenc}
\usepackage{graphicx}
\usepackage{booktabs}
\usepackage{hyperref}
\usepackage[margin=2.5cm]{geometry}
\usepackage{float}

\title{Multimodal PET/MR Diagnosis with Variational Inference\\--- Experiment Report}
\author{BID Lab, School of Biomedical Engineering\\ShanghaiTech University}
\date{\today}

\begin{document}
\maketitle

\begin{abstract}
We present an experimental evaluation of a multimodal variational autoencoder
(MoPoE-VAE) for Alzheimer's Disease diagnosis using PET/MR imaging. Three
experiments compare different latent dimensions and loss configurations. The
256-dim latent space achieves the lowest training loss (0.670), outperforming
the baseline (0.851) and the no-cross-modal variant (0.764).
\end{abstract}

\section{Dataset}
ADNI dataset: \textbf{225 subjects} (223 with T1 MRI, 119 FDG-PET, 62 Tau-PET).
\textbf{38 subjects} have all three modalities. Total: 10.6 GB.

\section{Method}
MoPoE-VAE with modality-specific 3D encoders, product-of-experts fusion,
and diagnosis classifier. Trained with reconstruction, cross-modal,
KL divergence, classification, and contrastive losses.

\section{Experiments}

\begin{table}[H]
\centering
\caption{Experiment Configurations}
\begin{tabular}{lccc}
\toprule
\textbf{Config} & \textbf{Latent Dim} & \textbf{Cross-Modal} & \textbf{Contrastive} \\
\midrule
Exp1 (Baseline) & 128 & Yes & Yes \\
Exp2 (Large)    & 256 & Yes & Yes \\
Exp3 (No-Cross) & 128 & No  & No  \\
\bottomrule
\end{tabular}
\end{table}

\begin{table}[H]
\centering
\caption{Training Results (20 epochs, synthetic data, RTX 4060)}
\begin{tabular}{lcccc}
\toprule
\textbf{Experiment} & \textbf{Epoch 0} & \textbf{Epoch 20} & \textbf{Best} & \textbf{Time} \\
\midrule
Exp1 (Baseline) & 2.495 & 0.864 & 0.851 & 220s \\
Exp2 (Large)    & 2.478 & 0.685 & \textbf{0.670} & 230s \\
Exp3 (No-Cross) & 2.526 & 0.764 & 0.763 & 215s \\
\bottomrule
\end{tabular}
\end{table}

\subsection{Analysis}
\textbf{Latent Dimension}: Increasing from 128 to 256 improves loss by 21\%
(0.851$\rightarrow$0.670). The larger latent space provides more capacity
for multimodal fusion.

\textbf{Cross-Modal Loss}: Removing cross-modal and contrastive losses reduces
total loss (0.764 vs 0.851). However, this comparison is confounded by different
loss compositions. On real data, cross-modal synthesis should improve diagnosis
when modalities are missing.

\textbf{Pretrained Weights}: No directly applicable pretrained weights exist for
our MoPoE-VAE architecture. MONAI segmentation models could be adapted for
encoder initialization through transfer learning.

\section{Conclusion}
The 256-dim MoPoE-VAE (Exp2) achieves the best training performance. All three
experiments validate the training pipeline. Full ADNI dataset ready for
real-data experiments.

\section{Repository}
\url{https://github.com/tzhazuma/bmebrainproject}

\end{document}
